% !TeX root = ../main.tex
\textbf{Влияние амплитудных и фазовых искажений (КИХ-дециматор)}

В завершающей части седьмого пункта было проанализировано влияние квадратурного дисбаланса цифрового гетеродина (амплитудной и фазовой асимметрии) на динамический диапазон системы при использовании единого КИХ-дециматора. Для моделирования применялся параметр \texttt{nco\_type = 'single'}.

Результаты моделирования в виде зависимостей SFDR от амплитудного (\texttt{amv}) и фазового (\texttt{pmv}) искажений представлены на рисунках \ref{fig:p7_sfdr_amv} и \ref{fig:p7_sfdr_pmv} соответственно.

\begin{figure}[H]
    \centering
    \includegraphics[width=0.8\textwidth]{SFDR_vs_Amplitude_Mismatch.png}
    \caption{Зависимость SFDR от амплитудных искажений при использовании КИХ-дециматора}
    \label{fig:p7_sfdr_amv}
\end{figure}

\begin{figure}[H]
    \centering
    \includegraphics[width=0.8\textwidth]{SFDR_vs_Phase_Mismatch.png}
    \caption{Зависимость SFDR от фазовых искажений при использовании КИХ-дециматора}
    \label{fig:p7_sfdr_pmv}
\end{figure}

\textbf{Сравнительный анализ и выводы:}
При идеальной настройке гетеродина (отсутствие искажений) архитектура с КИХ-дециматором обеспечивает выигрыш в динамическом диапазоне: начальное значение SFDR составляет $\approx 83$ дБн, что на 12 дБ лучше, чем у каскада с CIC-фильтром ($\approx 71$ дБн). Это преимущество достигается за счет лучшего подавления высокочастотных шумов квантования.

Однако при появлении малейшего дисбаланса наблюдается резкая деградация сигнала. Сравнение текущих графиков с графиками из пункта 4 выявляет важнейшую закономерность: \textbf{абсолютные значения деградировавшего SFDR при наличии искажений полностью совпадают для обеих архитектур фильтрации}. Например, при фазовой ошибке в $2^\circ$ уровень SFDR падает до $\approx 35$ дБн как при использовании CIC-фильтра, так и при использовании КИХ-дециматора. Аналогичная картина наблюдается и для амплитудного дисбаланса: при \texttt{amv} = $0.2$ дБ SFDR в обоих случаях опускается до $\approx 39$ дБн.

Физическая и математическая причина такого поведения заключается в природе зеркальной помехи:
\begin{itemize}
    \item В отличие от шумов квантования, которые распределены по широкой полосе частот и могут быть отфильтрованы, зеркальная помеха, возникающая из-за неортогональности I и Q каналов, попадает \textit{строго внутрь рабочей полосы пропускания} фильтра.
    \item Мощность этой помехи определяется исключительно тригонометрическими соотношениями внутри перемножителей гетеродина и никак не зависит от параметров последующего фильтра нижних частот.
    \item Цифровой фильтр (будь то CIC или КИХ) пропускает эту помеху на выход без изменений, так как она является для него низкочастотным сигналом.
\end{itemize}

Следовательно, превосходные характеристики подавления КИХ-дециматора в полосе задерживания оказываются абсолютно бесполезными против зеркальных помех. Это доказывает, что в системах квадратурной цифровой обработки главным лимитирующим фактором качества является математическая точность формирования квадратурных колебаний, а не выбор архитектуры децимирующего фильтра.